\chapter{Extended Worked Examples}
\label{app:extended-examples}

\begin{quotation}
\noindent\textit{Synopsis.} This appendix sketches additional domains ---
textual transmission, oral tradition, and biological inheritance ---
that further test and illustrate the framework beyond the chapters given
full treatment in Part~\ref{part:worked-domains}. Treatment here is at
appendix depth: each example is developed far enough to show the fit is
real, not far enough to constitute a chapter.

\end{quotation}

\section{Manuscript Transmission}

A chain of scribal copies $A_0$ (the autograph) $\to A_1 \to \cdots \to
A_n$ (an extant manuscript) is a substrate chain in the sense of
\cref{def:substrate-chain}, with each transition governed by a
copyist's own budget: time, eyesight, familiarity with the source
language, and attentiveness all bound $\beta_i$. Unlike
Chapter~\ref{ch:maps}'s scale bar, a scribal error is essentially never
declared: a miscopied word is presented exactly as confidently as a
correctly copied one, making most manuscript corruption an instance of
\cref{def:corrupting-loss} rather than \cref{def:benign-loss}.

The chain structure is not merely formally convenient here; it is the
literal basis of the scholarly discipline of textual criticism, which
reconstructs a stemma --- a family tree of copies --- from surviving
manuscripts and infers readings of the lost autograph by comparing
independent branches.

\begin{proposition}[Branch Agreement as Evidence for the Archetype]
\label{prop:stemma-agreement}
If two branches of a manuscript stemma, descended independently from a
common ancestor, agree on a reading, and copying errors are independent
across branches, that reading is more likely to reflect the ancestor
than a reading found in only one branch.
\end{proposition}

\begin{proofsketch}
This is \cref{def:branch-agreement} applied to the manuscript domain: if
errors in branch $X$ and branch $Y$ are independent, the probability
that both branches independently introduce the \emph{same} erroneous
reading is far lower than the probability that both correctly preserve a
shared true one. Agreement is therefore stronger evidence for the
ancestor's reading than either branch's testimony alone. This is the
formal content behind the standard editorial practice of preferring
readings attested across independent branches of a stemma.
\end{proofsketch}

Where branches disagree, an editor constructing a critical text must
choose a reading without being able to fully recover which branch, if
either, preserved the ancestor correctly -- a decision structurally
similar to \cref{sec:coord-adjudication}'s treatment of adjudication: the
editor's judgment substitutes an operative reading for an underlying
fact that may never be recoverable, though editorial judgment aims at
recovering the truth in a way judicial adjudication does not always
claim to.

\section{Oral Tradition}

Each retelling of an oral epic is a transition in the sense of
\cref{prop:memory-relay}: a performer's schema, operating on whatever
residual trace of prior tellings they carry, produces a new performance.
Oral-formulaic composition --- fixed epithets, metrical formulae,
repeated type-scenes, of the kind documented in the study of oral epic
tradition --- functions as a strong negative-interference device in the
sense of \cref{def:interference}: because performer and audience share
the same formulaic expectations, maintaining a formula costs the
performer far less than reconstructing an equivalent passage from
scratch would, exactly as \cref{rem:priors-reduce-cost} showed for a
painting's viewer.

This case surfaces a complication genuinely sharper than anything Part
IV's chapters needed to confront. Chapter~\ref{ch:budget-relay}'s root
object (\cref{def:root-object}) presupposes some $A_0$ existing prior to
the chain of transitions that follow it. An oral tradition without a
fixed textual original may have no such $A_0$ at all: each performance is
arguably a new production rather than a rendering of a prior fixed
telling, in which case the entire preservation/reconstitution question of
Part~\ref{part:modes-of-access} may not have a privileged answer to
converge toward, since there may be no single original for any
performance's fidelity (\cref{def:fidelity}) to be measured against. This
monograph does not resolve this; it is flagged here as a further, distinct
complication from the MEM|8 fork of Chapter~\ref{ch:mem8}, arising even
under the narrow reading, since it concerns whether a root object exists
at all rather than how faithfully a consumer accesses one. Where a
tradition does have something root-object-like -- a single elder who
knows a telling no one else does, as opposed to a story many performers
independently carry -- \cref{rem:witnesses-persistence-measure}
(Chapter~\ref{ch:four-operations}) offers unverified but concrete
vocabulary for exactly this distinction: the elder's version has
reconstruction redundancy $\rho_W = 1$, a single point of failure, while
the widely-carried story has $\rho_W \gg 1$.

\section{Biological Inheritance}

\begin{quotation}
\noindent\emph{A caveat before this section.} \emph{Conservation of
Ambiguity} reportedly applies repair-budget theory to biological
evolution directly. The sketch below has not been checked against that
treatment, and it is stated as an independent, appendix-level
observation rather than as an extension building on it; a careful
development of this case should verify the relationship before treating
either as complete.
\end{quotation}

Reproduction with recombination and mutation can be modeled as a
substrate chain: a parental genome $A_0$, an offspring genome $A_1$
produced by a transition governed by the fidelity of DNA repair and
replication machinery --- a budget in a sense closer to literal than
metaphorical, since the relevant biological machinery is itself called
repair machinery. Traits neutral under current selective pressure are
exactly the non-required distinctions of \cref{def:budget-discard}: they
are not funded preferentially, and under relaxed selection or genetic
drift they can be lost across generations, exactly as
\cref{prop:monotonicity} describes cumulative discard accumulating along
an uncoordinated chain.

\begin{proposition}[Environmental Change as Task Discovery]
\label{prop:environment-as-task}
If a trait is not required under the selective pressure $\tau_i$
operative during generations $0, \ldots, n-1$, and a later environmental
change constitutes a new selective task $\tau^*$ requiring that trait,
the trait's loss under $\tau_0, \ldots, \tau_{n-1}$ is unrecoverable by
\cref{prop:monotonicity}, regardless of how strongly $\tau^*$ would now
favor it.
\end{proposition}

\begin{proofsketch}
Immediate from treating successive generations as relay stages and
environmental change as introducing an unanticipated task, exactly in
the sense of \cref{prop:task-unavailability}. This is the formal
counterpart of a familiar evolutionary observation: a lineage cannot
retroactively retain a trait it has already lost to drift or relaxed
selection merely because a later environment would have favored it.
\end{proofsketch}

This proposition is offered cautiously, as a structural observation
about the fit between the two frameworks rather than as a contribution
to evolutionary theory, and its relationship to \emph{Conservation of
Ambiguity}'s own treatment of the same domain remains, per the caveat
above, to be checked.
